% SHARD-ID: RC-03-WHAT-RH-HAS-BECOME
% SHARD-TITLE: What the Riemann Hypothesis Has Become
% SHARD-SUMMARY: Translates the Riemann hypothesis through the one determinant identity into four equivalent costumes: counting fluctuations, subleading moduli of a small matrix, the Ramanujan property of a graph, and one correlation length of a transfer semigroup.
% SHARD-SUMMARY: Separates what the translation settles (the semigroup exists, the zeros are its modes, the functional equation is the pairing of its rates) from what it does not (a mechanism forcing each pair to be degenerate), and records the erratum on the words "entanglement Hamiltonian".
% SHARD-KEYWORDS: one identity, side A, side B, small matrix, Ramanujan, correlation length, explicit formula, pairing, Hilbert-Polya, erratum
% SHARD-NOTES: notes/what-rh-has-become.md
% SHARD-SCRIPTS: none

\section{What the Riemann Hypothesis Has Become}
\label{sec:what-rh}

Nothing here is new mathematics; it is a translation, written for TJO after the
founding session from the ingredients built up there. Its point is one sentence,
isolated in \cref{obs:rh-small-matrix} and repeated in every later costume: the
Riemann hypothesis says certain counts fluctuate no more than a random sequence would.

\subsection{The one identity, and its two readings}

\begin{observation}[The one identity]
\label{obs:one-identity}
\claimstatus{obs:one-identity}{sketched}
For any matrix $M$ with eigenvalues $\lambda_1,\dots,\lambda_N$,
\[
\frac{1}{\det(1-\uz M)} \;=\; \prod_i\frac1{1-\lambda_i\uz}
\;=\; \exp\Big(\sum_{\wl\ge1}\frac{\Tr(M^{\wl})}{\wl}\,\uz^{\wl}\Big).
\]
The middle expression is a bosonic partition function, one mode per eigenvector with
Boltzmann factor $\lambda_i\uz$; the right expression is a sum over closed walks of
length $\wl$, carrying the $1/\wl$ of a ring's rotation symmetry. One object, two
descriptions: a gas of modes and a gas of rings
(\cref{def:side-a-side-b,def:one-identity}).
\end{observation}

The whole subject is the tension between the two sides. For a permutation matrix the
ring side is concrete: $\Tr(M^{\wl})$ counts the fixed points of the $\wl$-th power,
the closed walks are the cycles and their repeats, and the identity
collapses to $\det(1-\uz M)^{-1} = \prod_{c}(1-\uz^{|c|})^{-1}$ over cycles $c$: a
Bose gas with one mode per cycle, energy the cycle length. That product is the Euler
product, and the primes will be the cycles.

\subsection{A finite permutation: the analogue is free}

\begin{observation}[RH for a finite permutation]
\label{obs:rh-finite-permutation}
\claimstatus{obs:rh-finite-permutation}{sketched}
A permutation matrix permutes an orthonormal basis, hence is unitary, hence has every
eigenvalue on the unit circle. The RH analogue at this level costs nothing.
\end{observation}

\subsection{An infinite permutation: a statement about a different matrix}

For infinitely many cycles the permutation matrix is still unitary and still useless,
its spectrum being roots of unity with infinite multiplicity. The count
$\operatorname{fix}(\wl) = \sum_{d\mid\wl}d\,c_d$ of fixed points of the $\wl$-th
power, with $c_d$ the number of $d$-cycles, remains meaningful, and the
interesting object is a \emph{small matrix} (\cref{def:small-matrix}) reproducing
them.

\begin{observation}[RH for the small matrix]
\label{obs:rh-small-matrix}
\claimstatus{obs:rh-small-matrix}{sketched}
A small matrix $M_{\mathrm{small}}$ with $\Tr M_{\mathrm{small}}^{\wl} =
\operatorname{fix}(\wl)$ for all $\wl$ exists exactly when the zeta function is
rational, and then its eigenvalues are new numbers, off the unit circle. Writing
$\operatorname{fix}(\wl) = \lambda_{\max}^{\wl} + \sum_{i\ge2}\lambda_i^{\wl}$,
the RH analogue is that every other nonzero $\lambda_i$ has $|\lambda_i| =
\sqrt{\lambda_{\max}}$ (\cref{def:rh-analogue}), so that the fluctuation has size
$\sqrt{\lambda_{\max}^{\wl}}$: the square root of the leading term, which is the
fluctuation of a sum of $\lambda_{\max}^{\wl}$ independent coin tosses. In counting
language, RH says the cycle counts fluctuate no more than a random sequence would.
\end{observation}

So the RH analogue is here a property some systems have and others do not. The note
records two shift examples, both perfectly good infinite permutations.

\begin{center}
\begin{tabular}{@{}lll@{}}
\toprule
shift & eigenvalues of the small matrix & RH analogue \\
\midrule
$4$-symbol cyclic rule & $2,\ 1+i,\ 1-i,\ 0$ & holds \\
golden-mean shift & $\phi,\ -1/\phi$ & fails \\
\bottomrule
\end{tabular}
\end{center}

\subsection{A graph: Ramanujan}

\begin{observation}[RH for a graph is the Ramanujan property]
\label{obs:rh-graph-ramanujan}
\claimstatus{obs:rh-graph-ramanujan}{sketched}
For a $(\qs+1)$-regular graph the permutation is the shift on non-backtracking closed
walks, the small matrix is the adjacency matrix $\Aadj$, and the Ihara--Bass formula
(\cref{def:ihara-bass}) compresses one to the other: each eigenvalue $\lambda$ of
$\Aadj$ produces a pair $\edgeev,\edgeev'$ of eigenvalues of the edge operator $\Hash$
through
\[
\edgeev^2 - \lambda\edgeev + \qs = 0,\qquad\text{so}\qquad
\edgeev\edgeev' = \qs .
\]
The leading $\lambda = \qs+1$ gives $\edgeev = \qs$, $\edgeev' = 1$. For the rest,
$\lambda^2 < 4\qs$ forces a complex conjugate pair with $|\edgeev| = |\edgeev'| =
\sqrt{\qs}$, on the critical circle, while $\lambda^2 > 4\qs$ gives a real pair, one
factor above $\sqrt{\qs}$ and one below. Hence RH for $\zetaX$ holds if and only if
$|\lambda|\le2\sqrt{\qs}$ for every nontrivial $\lambda$, i.e.\ iff the graph is
Ramanujan (\cref{def:ramanujan-graph}), and the closed non-backtracking walk counts
are $\Ncount{\wl} = \qs^{\wl} + O(\qs^{\wl/2})$.
\end{observation}

Three remarks carry over to the primes. First, \textbf{the pairing comes from a
symmetry, the hypothesis from a bound}: $\edgeev\edgeev' = \qs$ is a functional
equation pairing eigenvalues symmetrically about $\sqrt{\qs}$ on a logarithmic scale,
and RH is the extra statement that each pair is degenerate in modulus, which happens
exactly when it is a conjugate pair. Second, \textbf{self-adjointness is not enough}:
$\Aadj$ is symmetric, so $\lambda$ is real, which guarantees the pairing but not the
bound; the prism graph $C_{16}\times K_2$ is symmetric and fails it. Third,
\textbf{the bound is geometric}: $[-2\sqrt{\qs},2\sqrt{\qs}\,]$ is the spectrum of the
infinite tree covering the graph, so RH says the graph has no eigenvalue outside the
spectrum of its universal cover; the word is tempered.

\subsection{A matrix product state: one correlation length}

Let $\Tmat = \Sig$ be the transfer matrix of a matrix product state
(\cref{def:transfer-matrix}), so that $\Tr\Tmat^{n} = \ringnorm{n}$ is the ring norm
of the periodic state on $n$ sites (\cref{def:ring-norm}) and the zeta function of
\cref{obs:one-identity} is the generating function of ring norms. In canonical form
(\cref{def:canonical-form}) the leading eigenvalue is $1$ and the rest lie inside the
unit disc, $\lambda_k = e^{-1/\corrlen{k}+i\phase{k}}$
(\cref{def:correlation-length}), so $\ringnorm{n} = 1 + \sum_k
e^{-n/\corrlen{k}}e^{in\phase{k}}$.

\begin{observation}[RH for an MPS]
\label{obs:rh-mps-correlation-length}
\claimstatus{obs:rh-mps-correlation-length}{sketched}
Translating \cref{obs:rh-graph-ramanujan}, where the leading growth is $\qs$ per step
and the subleading modulus $\sqrt{\qs}$: in normalised units the subleading modulus is
$\qs^{-1/2}$, so $e^{-1/\corrlen{k}} = \qs^{-1/2}$ and $\corrlen{k} = 2/\log\qs =
2/\ent$ for every $k$. The RH analogue for an MPS reads: all correlation lengths are
equal, and equal to twice the inverse entropy rate. The phases $\phase{k}$ are
unconstrained and there may be infinitely many of them; only the moduli are pinned.
\end{observation}

The hierarchy this sits in has three levels, and each has a classical number-theoretic
counterpart, anticipating \cref{obs:rh-flow-pairing} below.

\begin{center}
\begin{tabular}{@{}p{3.1cm}p{2.5cm}p{4.0cm}p{4.7cm}@{}}
\toprule
statement about $\Tmat$ (or $\flowgen$) & name & counting meaning & number theory \\
\midrule
a unique eigenvalue on the unit circle &
primitivity (Perron--Frobenius) &
ring norms follow the leading term &
$\zeta(1+it)\ne0$, the prime number theorem \\
all others inside a smaller circle &
spectral gap &
fluctuation exponentially smaller than the leading term &
the zero-free region, $\cheb(x) = x + O(x e^{-c\sqrt{\log x}})$ \\
all others on \emph{one} circle of radius $e^{-1/\corrlen{k}}$ with
$\corrlen{k} = 2/\ent$ &
Ramanujan / RH &
fluctuation is the square root of the leading term &
all zeros on $\re s = \tfrac12$ \\
\bottomrule
\end{tabular}
\end{center}

\subsection{The primes: continuous length, and the sign}

For the primes the cycle lengths are $\log\pp$, rationally independent, so the matrix
powers become a semigroup $e^{\ringlen\flowgen}$ in a continuous length $\ringlen$ and
the ring norms become a measure in $\ringlen$. Start from the Riemann--von Mangoldt
formula for $\cheb$, valid for $x$ not a prime power, put $x = e^{\ringlen}$ and
differentiate in $\ringlen$: the left side becomes a sum of delta functions, one per
prime power, of weight $\log\pp$, and the right side a sum of exponentials.

\begin{observation}[The explicit formula as a trace, and the pairing]
\label{obs:rh-flow-pairing}
\claimstatus{obs:rh-flow-pairing}{sketched}
With $\cheb(x) = x - \sum_{\zero}x^{\zero}/\zero - \log2\pi - \tfrac12\log(1-x^{-2})$,
\[
\sum_n \vM(n)\,\delta(\ringlen-\log n) \;=\; e^{\ringlen}
\;-\; \sum_{\zero}e^{\zero\ringlen}
\;-\; \big(e^{-2\ringlen}+e^{-4\ringlen}+\cdots\big),
\]
which is $\Tr e^{\ringlen\flowgen}$ for a generator $\flowgen$ whose eigenvalues are
$1$ (the pole of $\zeta$, the leading term, entropy $1$), every zero $\zero$, and the
trivial zeros $-2,-4,\dots$: the trace identity of \cref{obs:one-identity} with $M^{\wl}$
replaced by $e^{\ringlen\flowgen}$, ring side the prime powers at lengths $\log n$,
mode side the zeros. The functional equation $\cxi(s) = \cxi(1-s)$ makes $\zero$ and
$1-\zero$ both zeros, so after dividing by the leading $e^{\ringlen}$ the normalised
decay rates come in pairs $(1-\sigs,\sigs)$ summing to $1$, exactly as
$\edgeev\edgeev' = \qs$ paired the graph eigenvalues around $\sqrt{\qs}$. That pairing
is a theorem; RH is the extra statement that every pair is degenerate, $\sigs =
1-\sigs = \tfrac12$.
\end{observation}

Two features of this trace identity are visible immediately. The first is \emph{the
sign}: the zeros enter with a minus sign relative to the leading term, the same minus
sign that marks the fermionic $H^1$ block in the elliptic-curve case
(\cref{sec:deligne}) and the absorption spectrum in Connes's language. The total
measure stays positive because $\vM(n)\ge0$, but the subleading sector of the bond
space contributes negatively, so whatever state this is, its bond space is graded, not
plainly bosonic. The second is the pairing just stated. The numerical check against
several thousand zeros is recorded in \cref{sec:riemann-channel}. Collecting the
costumes:

\begin{center}
\begin{tabular}{@{}p{2.4cm}p{11.6cm}@{}}
\toprule
language & the Riemann hypothesis says \\
\midrule
counting & $\cheb(x) = x + O(x^{1/2+\epsilon})$: the prime powers fluctuate like coin tosses \\
small matrix & every subleading mode has modulus the square root of the leading one \\
MPS & the state whose rings are the primes has \emph{one} correlation length, $\corrlen{k} = 2$ since $\ent = 1$, and an infinite tower of frequencies $\gamma_n$ \\
channel & the transfer semigroup $e^{\ringlen\flowgen}$ has a single uniform decay rate $\tfrac12$ \\
\bottomrule
\end{tabular}
\end{center}

\subsection{The Hilbert space where this is literally true}

The channel version is not only an analogy. It is a theorem on a concrete Hilbert
space: the model space $\KS$ of the scattering symbol $\Ssym$ of the modular surface,
on which the compressed dilation semigroup $\Zt$ has one mode per zero, with decay
rate and frequency read off from that zero. This is Lax--Phillips
\cite{laxphillips1976}, and ``RH iff every mode of $\Zt$ decays at the same rate'' is
Faddeev--Pavlov \cite{faddeev1972scattering}. What the founding session added is where
$\Ssym$ comes from: its phase is the imaginary part of the Bost--Connes Lévy exponent
at the critical temperature $\invtemp = 1$ \cite{bostconnes1995}, so the ring gas and
the transfer semigroup are one object seen from its two ends, bridged by the critical
phase. Construction, normalisations (the factor $\tfrac12$ from the doubled variable)
and numerics are in \cref{sec:riemann-channel};
\cref{def:scattering-symbol,def:inner-model-space,def:riemann-channel} fix the
objects.

\subsection{Why it is still hard, in this language}

The graph case is the honest guide. There side B exists, the trace identity holds and
the small matrix is symmetric for every graph, and most graphs fail the bound. Being
Ramanujan is never a consequence of the structure; it is an extra fact, and in every
known case it comes from arithmetic input of Deligne's type \cite{deligne1974weil}.
For Selberg's surfaces the same statement, no Laplace eigenvalue below $\tfrac14$, is
open.

What the translation settles: RH is a statement about the \emph{sizes} of the
subleading modes of a transfer semigroup, not about its existence. The semigroup
exists, the zeros are its modes, the primes are its rings, the functional equation is
the pairing of its rates. What it does not settle: a mechanism forcing the pairs to be
degenerate. In positivity form (\cref{def:weil-positivity}), RH holds iff $\Fweil$ is
positive-definite on the dilation group, iff its spectral measure is positive on the
real line, which by Bochner is the same as all its atoms being real.

\begin{observation}[Hermiticity of the Hilbert--P\'olya operator]
\label{obs:hilbert-polya-hermitian}
\claimstatus{obs:hilbert-polya-hermitian}{sketched}
Write the length-flow generator on the subleading sector as $\flowgen = -\tfrac12 +
iH$ (\cref{def:hilbert-polya}); then $H$ has eigenvalues $\gamma_n -
i(\sigs_n-\tfrac12)$, and RH says exactly that $H$ is Hermitian, i.e.\ that the
transfer semigroup is a uniform damping times an oscillation with real frequencies.
$H$ is the Hilbert--P\'olya operator. The non-Hermitian generator carries the complex
pairs and the Hermitian operator their symmetric combination, here $\zero(1-\zero) =
\sigs(1-\sigs) + \gamma^2 + i\gamma(1-2\sigs)$, which is real iff $\sigs = \tfrac12$;
$\flowgen$ is to $H$ what the Hashimoto edge operator $\Hash$ is to the adjacency
matrix $\Aadj$ of a graph. For graphs the Hermitian object is given and its reality is
a theorem; for the primes it is the conjecture.
\end{observation}

Proving it would mean exhibiting that state, or a Stinespring form of $\Zt$ with the
prime dilations as Kraus operators, in a way that makes the positivity manifest; that
is \cref{conj:stinespring-jump-form}. The graph lesson says the Kraus structure alone
will not do it, and that some arithmetic rigidity has to enter, playing the role that
weights play for curves (\cref{sec:deligne}).

\subsection{Erratum on words}

An earlier draft, and the founding session, called $-\log\Tmat$ the entanglement
Hamiltonian. That is wrong: the entanglement Hamiltonian of an MPS is $-\log$ of a
reduced density matrix built from the fixed points of $\Tmat$ alone, is Hermitian by
construction, and does not see the subleading spectrum. The operator whose spectrum
contains the zeros is the transfer generator, non-Hermitian in general, whose
eigenvalue phases are the oscillation frequencies of the ring norms.
\cref{obs:hilbert-polya-hermitian} states the correct relation; the route is listed
among the dead ones in \cref{sec:open}.

\subsection{Summary}

A zeta function is one partition function with two Fock-space descriptions: a gas of
rings (cycles, closed walks, prime powers) and a gas of modes (eigenvalues of a small
transfer matrix); the trace identity links them. For the primes the rings have lengths
$\log n$ and weights $\vM(n)$, the modes are the zeros of $\zeta$ plus the pole, and
the functional equation pairs the modes' decay rates as $(\sigs,1-\sigs)$. The Riemann
hypothesis says every pair is degenerate at $\tfrac12$: the transfer semigroup has one
correlation length, $\corrlen{k} = 2$, so the ring counts fluctuate by no more than
the square root of their size, so the primes are as regular as coin tosses. That
semigroup exists explicitly, its scattering symbol is the phase of the Bost--Connes
system at its critical temperature, and RH is the statement that this channel is
Ramanujan.
